

% ------------------------------------------------------------
\subsection{Request-Path Components}
\label{subsec:request-path}
% ------------------------------------------------------------

This subsection describes the synchronous, user-facing pipeline: the
sequence of components that a single user message---or a single image
request---traverses between arrival at the orchestrator and the
return of the rendered response. Asynchronous concerns such as event
logging, analytics rollups, and the advertiser portal are deferred to
Section~\ref{subsec:offline-path}. The components below are presented
in the order in which they execute on the request path, and each is
characterised by its inputs, its outputs, and the failure mode it
adopts when a downstream dependency is unavailable.

% ------------------------------------------------------------
\subsubsection{Orchestrator API (FastAPI)}
\label{subsubsec:orchestrator}
% ------------------------------------------------------------

The orchestrator is a single FastAPI service that exposes the public
HTTP surface of the platform and coordinates every internal module
involved in producing a response. It is deliberately monolithic:
context extraction, ad retrieval, response synthesis, session
management, and event logging are mounted as in-process modules behind
a thin router layer rather than as independent services. This choice
trades the horizontal scalability of a microservice topology for two
properties that matter more at the MVP scale, namely a single
deployment artefact and the elimination of inter-service network hops
from the critical latency path.

The principal endpoint is \texttt{POST /chat}, whose handler implements
a seven-stage pipeline. The handler first loads the conversation state
for the supplied \texttt{session\_id} from Redis and, if the previous
turn surfaced an advertisement, dispatches an engagement-detection
call against the user's new message before any other work proceeds.
The user's turn is then persisted, and the message together with the
rolling history is passed to the context extraction module. If the
extracted \texttt{ad\_receptivity} field is anything other than
\texttt{none}, the handler issues an embedding call on the extracted
context, retrieves a re-ranked candidate set from the ad service,
generates per-impression tracking URLs, and invokes response synthesis
with the chosen ad bound into the prompt. If receptivity is
\texttt{none}---the safety branch---the same synthesis call is made
with an empty ad list, so that sensitive conversations follow an
identical code path but never carry sponsored content. The handler
finally appends the assistant turn to Redis, logs an impression event
if an ad was injected, and returns a response payload containing the
rendered text and, when applicable, the ad metadata required by the
frontend to render the \texttt{[Sponsored]} badge.

Two design decisions deserve emphasis. First, every call into the AI
modules is wrapped in a guarded \texttt{try/except} block whose
fallback is to continue the pipeline without the failed signal: a
context-extraction failure degrades the request to an unsponsored
response rather than to an error, and a response-synthesis failure is
the only condition under which the endpoint returns a non-2xx status.
The orchestrator therefore treats advertising as a strictly additive
overlay on a working chat service, never as a precondition for it.
Second, the secondary endpoints---\texttt{GET /track/click} for click
attribution, \texttt{GET /analytics/summary} for the dashboard, and
the \texttt{/portal/*} family for advertiser self-service---are mounted
on the same application but execute outside the latency-critical chat
path and are therefore free to perform synchronous database work that
the chat handler would have to defer.

% ------------------------------------------------------------
\subsubsection{Context Extraction Module}
\label{subsubsec:context-extraction}
% ------------------------------------------------------------

The context extraction module is the first AI call on the request
path and the component on which every subsequent ad-related decision
depends. It receives the user's current message together with the
rolling conversation history and returns a structured object whose
fields are the intent label, a list of topics, a list of named
entities, a coarse sentiment value, a continuous purchase signal in
the closed unit interval, and a four-valued ad-receptivity label
drawn from \(\{\texttt{high}, \texttt{medium}, \texttt{low}, \texttt{none}\}\).
This object is the sole interface between raw conversational text and
the retrieval and synthesis stages that follow; no downstream module
re-reads the user message for ad-related decisions.

The module is implemented as a single call to a lightweight chat
model, chosen so that the extraction step contributes only marginally
to the end-to-end latency budget while still benefiting from the
in-context reasoning that distinguishes, for example, a product
research query from a venting message that happens to mention a brand.
The prompt instructs the model to return strict JSON conforming to the
schema above, and the orchestrator validates the response against a
Pydantic model before propagating it; a parse failure is treated as a
soft error and degrades the request to the unsponsored branch.

The \texttt{ad\_receptivity} field carries a load that the other
fields do not. It is the platform's safety gate: when the extractor
labels a message as describing a health crisis, an episode of
emotional distress, or a financial hardship, the field is set to
\texttt{none} and the orchestrator unconditionally skips ad
retrieval and ad injection for that turn. This places the safety
decision at the earliest possible point in the pipeline, before any
embedding or retrieval work has been performed, and ensures that
sensitive conversations cannot incur an impression even
accidentally. The remaining three values are interpreted as a soft
gate by the re-ranker rather than as hard switches, so that a
borderline-receptive turn can still produce a recommendation when the
semantic match to inventory is sufficiently strong.

% ------------------------------------------------------------
\subsubsection{Session Management (Redis)}
\label{subsubsec:session-management}
% ------------------------------------------------------------

Session state is held in a single Redis instance and keyed by the
\texttt{session\_id} that the frontend generates at the start of each
conversation. Each session record contains a rolling window of the
ten most recent turns, an accumulated list of topics observed across
the conversation, the history of purchase signals reported by the
context extractor, the identifiers of advertisements already shown in
the session, and a turn counter. The rolling window bounds both the
memory footprint per session and the number of tokens that the
context extractor and the response synthesiser must process on each
turn, which is the dominant cost driver in both calls.

Three properties of this design are worth noting. First, the history
of purchase signals enables the re-ranker to detect monotone
escalation across turns, which is a stronger purchase indicator than
any single message in isolation; the platform exploits this without
requiring the LLMs themselves to reason over multi-turn dynamics.
Second, the set of advertisements already shown is consulted by the
re-ranker to suppress immediate repetition, which would otherwise be
the dominant failure mode of a purely similarity-driven retriever in
a long conversation about a single topic. Third, because Redis is the
only stateful dependency on the chat path that the orchestrator
writes to synchronously, the platform can tolerate a Redis outage by
falling back to a stateless mode in which every turn is treated as
the first turn of a fresh conversation; the user observes a loss of
personalisation but not a loss of service.

% ------------------------------------------------------------
\subsubsection{Response Synthesis with Ad Injection}
\label{subsubsec:response-synthesis}
% ------------------------------------------------------------

Response synthesis is the second and more expensive of the two LLM
calls on the request path. It receives the user's message, the
rolling conversation history, the structured context object produced
by the extractor, the chosen advertisement (if any), and the
pre-generated tracking URL associated with that impression. Its
output is a single markdown-formatted assistant message which the
frontend renders directly into the chat transcript.

The defining design decision is that ad injection is realised as
\emph{prompt-level conditioning} rather than as post-hoc text
rewriting. When an advertisement is selected, the synthesiser's
system prompt is augmented with a sponsored block that contains the
product name, its description, the advertiser-supplied creative copy,
any campaign-level presentation guidance, and a strict instruction to
weave the product into the response as a markdown link of the form
\texttt{[Sponsored] [\textit{product}](\textit{tracking\_url})},
followed by a single compelling reason. The synthesis model is then
free to determine where in its response the sponsored mention is
introduced, which substantive parts of the user's question to address
first, and how to phrase the recommendation so that it follows
naturally from the surrounding answer. This delegation is what
distinguishes the platform from a banner-style overlay: the model is
not asked to attach an advertisement to a response, it is asked to
write a response in which the advertisement participates.

When the receptivity gate has fired and no advertisement is supplied,
the same synthesis call is made with an unmodified system prompt. The
two branches share a single code path and a single model invocation,
which guarantees that the unsponsored response distribution is not
distorted by structural differences between the two branches and that
a regression in the sponsored branch cannot silently corrupt the
unsponsored one. The \texttt{[Sponsored]} marker and the tracking URL
are the only artefacts by which a downstream observer---the user, the
analytics dashboard, or a future audit---can distinguish a sponsored
response from an organic one.

% ------------------------------------------------------------
\subsubsection{Image Generation with Product Placement}
\label{subsubsec:image-generation}
% ------------------------------------------------------------

The image generation endpoint is the platform's second generative
surface and the visual analogue of the chat pipeline described above.
It is exposed as a distinct route on the same orchestrator and
operates under a deliberately looser latency budget, since
diffusion-model inference routinely exceeds ten seconds and is
dominated by GPU time rather than by the orchestration overhead that
constrains the text path. Its purpose is to demonstrate that the
context-extraction and retrieval-augmented matching pipeline
generalises beyond text and can deliver organic, RAG-matched product
placement directly into user-requested imagery.

When a user submits an image prompt, the orchestrator routes the
prompt through the same context extractor used on the chat path,
issues an embedding call against the resulting context object, and
retrieves a re-ranked candidate set from the ad service. If a
sponsored product is selected, the original prompt is not passed to
the image model directly. Instead, an auxiliary lightweight LLM call
rewrites the prompt into a \emph{placement prompt}: a new
specification that preserves the user's stated subject, style, and
composition while introducing the matched product as a diegetic
element of the scene---worn by a subject, held by a hand, resting on
a surface, displayed on a shelf---chosen to fit the semantics of the
original request rather than to dominate it. The rewritten prompt is
then passed to the image model, \texttt{gpt-image-1}, which renders
the final image.

A second branch of the pipeline activates when the matched
advertisement carries a reference photograph of the product in the ad
inventory. In this case the rewritten prompt is dispatched to the
model's image-edit endpoint together with the reference image, so
that the rendered product is visually grounded in the actual article
rather than reconstructed from the model's prior over the product
name. This is a small change in the API call but a significant change
in the semantics of the output: it is what allows a generated scene
to depict a recognisable, brand-accurate product rather than a
plausible-looking generic substitute, and it is what makes the visual
pipeline a meaningful extension of the textual ``instruct the
generative model to weave in the ad'' design philosophy rather than a
cosmetic copy of it. As with the chat path, a failure anywhere in the
ad-related branches degrades the request to an unsponsored image
generation: the user always receives an image, and the sponsored
overlay is strictly additive.